\section{Worum es geht}

\textbf{[SETZUNG]} $\xi = 4/30000$ ist eine Setzung (A020). Dieses Dokument
beantwortet die Frage, ob diese Zahl \emph{nur} empirisch verankert ist oder ob
sie aus der Strukturdes Systems folgt. Das Ergebnis: $\xi$ ist nicht beliebig.
Der Exponent $-4$ in der Größenordnung $10^{-4}$ folgt aus der fraktalen
Raumzeitdimensionalität $D_f=3-\xi$ und der Vier-Dimensionalität des Universums.
Der Massenskalierungsexponent $\kappa=7$ ist die \emph{einzige selbstkonsistente
Lösung} des e-p-$\mu$-Dreiecks.
\textit{(Vermerk vom 1.~Okt.~2026: $\kappa=7$ ist eine Wahl der Zerlegung
$R_{pe} = K\cdot(4/3)^7$ mit Restterm $K$, keine Selektion; siehe den Vermerk im
Abschnitt zum e-p-$\mu$-System.)}
\textit{(Entfernt am 1.~Okt.~2026: vorher \glqq Der Vorfaktor $4/3$ ist durch den
Casimir-Effekt unabhängig bestätigt\grqq{} und im Folgesatz \glqq dreifach
verankert\grqq{} -- die Casimir-Energie pro Fläche ist
$-\pi^2\hbar c/(720\,a^3)$ ohne Faktor $4/3$; eine Casimir-Bestätigung des
Faktors $4/3$ gibt es nicht.)}
Damit ist die Größenordnung von $\xi$ verankert
-- aber \emph{nicht} aus den drei Setzungen von A020 allein ableitbar; ein
gemessener Massenanker bleibt nötig.

\section{Das Zirkularitätsproblem}

\textbf{[S]} Die ursprüngliche Darstellung scheint zirkulär:
\[
  \frac{m_p}{m_e} = 245\times\Bigl(\frac43\Bigr)^7
  \;\Rightarrow\; \xi = \frac{4}{30000}.
\]
\textbf{Einwand:} Warum gerade $\kappa=7$? Warum $K=245$? Ist das nicht
Rückwärts-Fitting?

\textbf{[B]} Die Antwort: der Exponent $\kappa=7$ wird nicht angepasst, sondern
er ist die \emph{einzige konsistente Lösung} für das vollständige e-p-$\mu$-System.
Das ist der Kern des Dokuments; der Rest zeigt es.
\textit{(Vermerk vom 1.~Okt.~2026: Die Aussage hält nicht; $\kappa=7$ ist eine
Wahl der Zerlegung, keine Selektion -- siehe den Vermerk im folgenden Abschnitt.)}

\section{Das e-p-$\mu$-System als Beweis}

\textbf{[K]} Die drei fundamentalen Massenverhältnisse:
\begin{align}
  R_{pe}    &= \frac{m_p}{m_e}   = 1836{,}15267343, \notag\\
  R_{\mu e} &= \frac{m_\mu}{m_e} = 206{,}7682830, \notag\\
  R_{p\mu}  &= \frac{m_p}{m_\mu} = 8{,}880. \notag
\end{align}
Die Multiplikativität erzwingt $R_{pe} = R_{\mu e}\times R_{p\mu}$. Für den
Ansatz $R_{pe} = K\times(4/3)^\kappa$ muss man nur den Exponenten testen:

\begin{center}
\renewcommand{\arraystretch}{1.3}
{\small
\begin{tabular}{rccr}
\toprule
$\kappa$ & Vorhersage $R_{pe}$ & Konsistenz & Fehler \\
\midrule
6 & $245\times(4/3)^6 = 1376{,}6$ & \texttimes & 25{,}0\,\% \\
\textbf{7} & $\mathbf{245\times(4/3)^7 = 1835{,}4}$ & \checkmark & \textbf{0{,}04\,\%} \\
8 & $245\times(4/3)^8 = 2447{,}2$ & \texttimes & 33{,}3\,\% \\
\bottomrule
\end{tabular}
}
\renewcommand{\arraystretch}{1.0}
\end{center}

\textbf{[K]} $\kappa=7$ ist die einzige konsistente Lösung -- kein Fitting,
sondern Selektion. Der Exponent folgt direkt:
\[
  \kappa = \frac{\ln(R_{pe}/K)}{\ln(4/3)}
  = \frac{\ln(1836{,}15/245)}{\ln(1{,}3333)} \approx 7{,}000.
\]
\textit{(Vermerk vom 1.~Okt.~2026: Die Tabelle hält das an $\kappa=7$ angepasste
$K=245$ fest ($K = R_{pe}/(4/3)^7$, Abschnitt zu $K$) und kann daher nur
$\kappa=7$ auszeichnen. Für jedes $\kappa$ passt $K_\kappa = R_{pe}/(4/3)^\kappa$
exakt: $\kappa = 6,\ 7,\ 8$ ergeben $K = 326{,}8;\ 245{,}1;\ 183{,}8$. Auch die
$\kappa$-Formel setzt $K = 245$ voraus. $\kappa=7$ ist eine Wahl der Zerlegung
$R_{pe} = K\cdot(4/3)^7$ mit Restterm $K$, keine Selektion; die Tabelle zeigt
keine Eindeutigkeit.)}

\section{Die fundamentale Begründung der $10^{-4}$-Skalierung}

\textbf{[B]} Warum gerade $10^{-4}$? Das Argument ist dreistufig:

\textbf{(1) Aus der fraktalen Dimension.} Die fraktale Dimension
$D_f=3-\xi\approx2{,}9998667$ erzeugt eine diskrete Skalenhierarchie. Die
relative Abweichung von der ganzzahligen Dimension:
\[
  \delta = 1-\frac{D_f}{3} = 1-\frac{3-\xi}{3} = \frac\xi3
  \;\approx\; 4{,}44\times10^{-5},
\]
also $\xi=3\delta=3-D_f$; $\xi$ ist die absolute Abweichung der fraktalen
Dimension von~3.
\textit{(Korrigiert am 1.~Okt.~2026: vorher $\delta \approx 1{,}333\times10^{-4}$
und \glqq Das macht $\xi$ zur relativen fraktalen Abweichung\grqq{} -- mit
$D_f = 3-\xi$ ist $\delta = \xi/3 = 4{,}444\times10^{-5}$; $1{,}333\times10^{-4}$
ist $\xi$ selbst, und die relative Abweichung ist $\delta$, nicht $\xi$.)}

\textbf{(2) Aus der Raumzeit-Dimensionalität.} In $d$-dimensionalen Räumen
erwartet man natürliche Skalierungen $\xi_d\sim(10^{-1})^d$. Für $d=4$:
\[
  \xi_4\sim(10^{-1})^4 = 10^{-4}.
\]
Das erklärt die Größenordnung; der Vorfaktor $4/3$ folgt daraus nicht (zu
Punkt (3) siehe dort).

\textbf{(3) Aus dem Casimir-Effekt.} Die Casimir-Energie pro Fläche zwischen
idealen parallelen Platten ist
\[
  \frac{E_\text{Casimir}}{A} = -\frac{\pi^2\hbar c}{720\,a^3}
\]
ohne Faktor $4/3$; eine Casimir-Bestätigung des Faktors $4/3$ gibt es nicht.
\textit{(Korrigiert am 1.~Okt.~2026: vorher \glqq Der Casimir-Effekt liefert den
Faktor $4/3$ unabhängig von Massenfits\grqq{} mit
$E_\text{Casimir} = -\pi^2\hbar c/(720\,a^3)\times 4/3$ sowie im vorigen Punkt
\glqq den genauen Vorfaktor $4/3$ liefert Punkt (3)\grqq{} -- Standardergebnis:
$E/A = -\pi^2\hbar c/(720\,a^3)$, Druck $\pi^2\hbar c/(240\,a^4)$,
Energiedichte $\pi^2\hbar c/(720\,a^4)$; ein Faktor $4/3$ kommt nicht vor und
wäre mit $+33\,\%$ durch die Messungen (etwa $1\,\%$ in der
Kugel-Platte-Geometrie) ausgeschlossen.)}
Nur die Quarte $(4/3)$ -- und nicht Quinte $(3/2)$ oder Terz $(5/4)$ -- liefert
konsistente Ergebnisse für das Massenspektrum. Die Konsistenz über fünf
unabhängige Verhältnisse:

\begin{center}
\renewcommand{\arraystretch}{1.3}
{\small
\begin{tabular}{lccc}
\toprule
Verhältnis & Experiment & T0 mit $\kappa=7$ & Fehler \\
\midrule
$m_p/m_e$       & 1836{,}1527 & 1835{,}4 & 0{,}04\,\% \\
$m_\mu/m_e$     & 206{,}7683  & Eingabe & -- \\
$m_p/m_\mu$     & 8{,}880     & 8{,}877 & 0{,}04\,\% \\
$m_\tau/m_\mu$  & 16{,}817    & Eingabe & -- \\
$m_n/m_p$       & 1{,}001378  & 1{,}001333 & 0{,}004\,\% \\
\bottomrule
\end{tabular}
}
\renewcommand{\arraystretch}{1.0}
\end{center}
\textit{(Korrigiert am 1.~Okt.~2026: vorher in der Spalte \glqq T0 mit
$\kappa=7$\grqq{} $m_\mu/m_e = 206{,}768$ ($0{,}001\,\%$),
$m_p/m_\mu = 8{,}880$ ($0{,}02\,\%$), $m_\tau/m_\mu = 16{,}817$
($0{,}02\,\%$) -- der Ansatz $K\cdot(4/3)^\kappa$ liefert nur
$m_p/m_e = 245\cdot(4/3)^7 = 1835{,}43$; daraus
$m_p/m_\mu = 1835{,}43/206{,}768 = 8{,}877$ ($-0{,}04\,\%$).
$m_\mu/m_e$ und $m_\tau/m_\mu$ standen als Messwerte ohne Rechenweg in der
Spalte und sind Eingaben; eine Konsistenz über fünf unabhängige Verhältnisse
belegt die Tabelle nicht.)}

\section{Was $K = 245$ ist und wofür man es braucht}

\textbf{[K]} $K=245$ ist ein \emph{Restterm}, kein fundamentales Objekt.
$\ln(R_{pe})/\ln(4/3) = 26{,}12$ -- nicht ganzzahlig. Um einen ganzzahligen
Exponenten $\kappa=7$ zu erhalten, zerlegt man:
\[
  K = \frac{R_{pe}}{(4/3)^7}
  = \frac{1836{,}152}{7{,}492} \approx 245{,}1.
\]
$K$ benennt den Rest nach dem Herausziehen von $(4/3)^7$. Es taucht in keiner
anderen FFGFT-Formel auf. Wäre $R_{pe}$ exakt $245\times(4/3)^7=1835{,}43$,
wäre $K$ exakt ganzzahlig; der kleine Rest $R_{pe}-245\times(4/3)^7=0{,}73$
erzeugt die 0{,}04\,\%-Abweichung.

\textbf{[S]} Die Primfaktorzerlegung $245=5\times7^2$ ist eine
post-hoc-Beobachtung -- $7$ erscheint wegen $\kappa=7$, $5$ ist Zufall.
Eine geometrische Ableitung von $K=245$ gibt es nicht; eine entsprechende
Formel ($245\approx\phi^{12}/(1-\xi)^2$) aus dem Altkorpus ist rechnerisch
falsch ($\phi^{12}/(1-\xi)^2\approx322$, Abweichung 31\,\%).

\textbf{[K]} \textbf{Ehrliche Einschätzung.} Die drei Verankerungen —
fraktale Dimension (1), Raumzeit-Dimensionalität (2), Casimir (3) — erklären
die Größenordnung $10^{-4}$ und den Vorfaktor $4/3$. Sie hängen aber alle am
e-p-$\mu$-System als empirischem Anker. Eine Ableitung von $\xi$ aus den drei
Setzungen von A020 allein (ohne gemessene Massen) gibt es nicht. Was vorliegt,
ist das stärkste verfügbare Argument: die Konsistenz der einzigen selbstkonsistenten
Lösung über fünf Massenverhältnisse.
\textit{(Vermerk vom 1.~Okt.~2026: Die Casimir-Verankerung (3) entfällt, siehe
dort; $\kappa=7$ ist eine Wahl der Zerlegung, keine Selektion, und die
Konsistenztabelle enthält zwei Eingaben -- siehe die Vermerke oben.)}

\section{Prüfskript}

\begin{itemize}[leftmargin=2em,itemsep=2pt,topsep=2pt]
  \item Numerische Gegenprobe: $\kappa=7$ als einzige konsistente Lösung für
    $R_{pe}$, $R_{\mu e}$, $R_{p\mu}$; Fehler für $\kappa=6,7,8$:
    \texttt{python/A\_Serie\_Skripte/a015\_xi\_ursprung.py}
    \textit{(Vermerk vom 1.~Okt.~2026: Ein Fehlervergleich bei festem
    $K=245$ zeigt keine Eindeutigkeit von $\kappa$, siehe den Vermerk zum
    e-p-$\mu$-System.)}
\end{itemize}

\section{Was offen bleibt}

\textbf{Ein Beweis von $\xi$ wird es nicht geben -- und das ist kein Mangel.}
$\xi$ ist ein Fundamentalparameter der Theorie -- analog zu $\alpha$ im
Standardmodell oder $G$ in der Gravitation. Fundamentalparameter werden nicht
aus tieferen Prinzipien bewiesen; sie werden durch Plausibilitätsgründe
motiviert, durch Konsistenz gestützt und durch Messungen bestätigt. Eine
„rein geometrische Ableitung ohne empirischen Eingang" ist für einen
Fundamentalparameter prinzipiell nicht möglich -- das gilt für jede Theorie.

\textbf{Was es gibt: plausible Gründe auf mehreren Ebenen.}

\textbf{(1) Die Quarte als einzige konsistente Basis.} Von allen harmonischen
Intervallen -- Quarte $4/3$, Quinte $3/2$, Terz $5/4$ -- trifft nur die
Quarte das e-p-$\mu$-System konsistent. Das ist kein Zufall der Tabelle,
sondern der Selektionsmechanismus des Experiments.
\textit{(Vermerk vom 1.~Okt.~2026: Mit freiem Restterm $K$ passt jede Basis und
jedes $\kappa$ exakt; die Auszeichnung von $4/3$ und $\kappa=7$ ist eine Wahl,
kein Selektionsmechanismus -- siehe den Vermerk zum e-p-$\mu$-System.)}

\textbf{(2) Fraktale Abweichung.} $\xi=3-D_f$ ist die absolute Abweichung
der fraktalen Dimension $D_f=3-\xi$ von der ganzzahligen Dimension~3.
\textit{(Korrigiert am 1.~Okt.~2026: vorher \glqq $\xi=1-D_f/3$ ist die relative
Abweichung\grqq{} -- $1-D_f/3 = \xi/3$, also $\xi = 3(1-D_f/3) = 3-D_f$.)}
Die Gleichung $\xi=\xi$ ist zirkulär, aber die Konsistenz ($D_f$ und $\xi$
bestimmen sich gegenseitig eindeutig) ist ein strukturelles Argument.

\textbf{(3) Natürliche Skalierung.} In $d$-dimensionalen Räumen ist $\xi_d
\sim(10^{-1})^d$ die natürliche Skalierung; für $d=4$ (3 Raum + 1 Zeit)
folgt $\xi_4\sim10^{-4}$.

\textbf{(4) Primfaktorstruktur.} $\xi=4/30000=2^2/(3\cdot2^4\cdot5^4)
=1/(3\cdot2^2\cdot5^4)$ -- Faktor~3 entspricht den Raumdimensionen,
Faktor~4 den Raumzeit-Dimensionen, $5^4$ der fraktalen Struktur.

\textbf{(5) Musiktheorieanalogon.} Pythagoreische 3:4:5-Tetraeder als
strukturelle Parallele zu den harmonischen Verhältnissen -- kein Beweis,
aber eine methodische Herkunft (Akkordeonbau, Akustik, Resonanzphysik),
die das Intervall 4:3 als natürliche Einheit nahelegt.

Diese Gründe zusammen motivieren $\xi=4/30000$ als plausiblen, konsistenten
Wert -- nicht als bewiesenen. Das ist der ehrliche Stand, und es ist derselbe
Stand wie bei jedem anderen Fundamentalparameter der Physik.

\textbf{$K=245$ ist Restterm, kein Fundamentalparameter.} Er folgt aus $R_{pe}$
und $\kappa=7$ -- ein Zwischenwert, keine eigenständige Größe (A015, Abschnitt
über K=245).

\section{Ersetzt}

Dieses Dokument nimmt auf: Dok.~009 ($\xi$-Ursprung: Zirkularitätsproblem,
e-p-$\mu$-System, Casimir-Bestätigung \textit{(Vermerk vom 1.~Okt.~2026: nicht
haltbar, siehe Punkt (3) der $10^{-4}$-Begründung)}, $K=245$, vollständiges Konsistenz-System),
Dok.~180 (vollständige Ableitungskette in fünf Schritten: T0-Lagrangian
$\to$ Selbstkonsistenzbedingung $\to$ Eindeutigkeit $\xi$ $\to$ $G$ $\to$ $L_0$,
Verbindung zum SI-System).

\section{Verwendung von KI-Werkzeugen}

Dieses Dokument ist unter Verwendung KI-gestützter Werkzeuge entstanden. Die
Fragestellungen, die Setzungen und alle inhaltlichen Entscheidungen stammen vom
Autor; die Werkzeuge formulieren aus, rechnen nach und erstellen Prüfskripte.
Die Verantwortung für Inhalt und Fehler liegt vollständig beim Autor. Der
ausführliche Wortlaut steht in A010.
